\documentclass{amsbook}

\newtheorem{theorem}{Theorem}[chapter]
\newtheorem{lemma}[theorem]{Lemma}

\theoremstyle{definition}
\newtheorem{definition}[theorem]{Definition}
\newtheorem{example}[theorem]{Example}
\newtheorem{xca}[theorem]{Exercise}

\theoremstyle{remark}
\newtheorem{remark}[theorem]{Remark}

\numberwithin{section}{chapter}
\numberwithin{equation}{chapter}

\begin{document}

\frontmatter
\title{Minqi's Solutions to the Exercises and Problems of Cormen, Leiserson, Rivest, Stein (2009), Introduction to Algorithms (3rd Edition)}

\author{Minqi Pan}

\date{\today}

\maketitle

\setcounter{page}{4}

\tableofcontents

\mainmatter

\chapter{The Role of Algorithms in Computing}

\section{Exercise 1.1-1}
Give a real-world example that requires sorting or a real-world example that requires computing a convex hull.
\begin{proof}[Answer]
Microsoft Windows provides a feature to sort all files in a folder by their sizes.

Computing a convex hull is required in ethology to estimate an animal's home range based on the points where the animal has been observed.
\end{proof}

\section{Exercise 1.1-2}
Other than speed, what other measures of efficiency might one use in a real-world setting?
\begin{proof}[Answer]
Number of CPU instructions, and memory usage.
\end{proof}

\section{Exercise 1.1-3}
Select a data structure that you have seen previously, and discuss its strengths and limitations.
\begin{proof}[Answer]
A red-black tree automatically keeps its height small in the face of arbitrary item insertions and deletions. Thus it has the strength of faster search operations which takes time proportional to the height of the tree. But it has the limitation of slowing down insertion and deletion operations because of additional rebalancing procedures. 
\end{proof}

\section{Exercise 1.1-4}
How are the shortest-path and traveling-salesman problems given above similar? How are they different?
\begin{proof}[Answer]
Both are trying to optimize path lengths to the shortest.

The traveling-salesman problem has more constraints than the shortest-path problem, in that it requires all nodes to be visited and the path to be a loop. The shortest-path problem, on the other hand, does not require any particular set of nodes to be visited other than the source node and/or the destination node, nor does it require the path to be a loop.
\end{proof}

\section{Exercise 1.1-5}
Come up with a real-world problem in which only the best solution will do. Then come up with one in which a solution that is ``approximately'' the best is good enough.
\begin{proof}[Answer]
Consider a delivery company with a central depot. Each day, it loads up each delivery truck at the depot and sends it around to deliver goods to several addresses. At the end of the day, each truck must end up back at the depot so that it is ready to be loaded for the next day.

Suppose that the company has already purchased many delivery trucks and is trying to determine whether the fuel tank of each truck is big enough to finish a one-day delivery tour without refueling. Also suppose that the fuel tank capacities are relatively small, so that it is not obvious to decide. Then only the best solution (i.e., the shortest tour route) can answer this question.

On the other hand, suppose that the delivery company has not purchased its delivery trucks yet, then the company can use the ``approximately'' best solution (i.e., the ``approximately'' shortest tour route) to decide the minimal fuel-tank capacity of the trucks to purchase.
\end{proof}

\section{Exercise 1.2-1}
Give an example of an application that requires algorithmic content at the application level, and discuss the function of the algorithms involved.
\begin{proof}[Answer]
The Nest Learning Thermostat is a smart thermostat that optimizes heating and cooling of homes and businesses to conserve energy. This application requires energy load management and HVAC control algorithms to make it smart and helpful, whether that be through sensing, machine perception, control systems, pattern recognition, or preference learning. The function of the algorithms involved is to keep its users comfortable while saving them money and reducing the world's energy consumption. 
\end{proof}
\end{document}
